\section{The original full packet at its formal boundary}
\label{GraphV2:sec:fpb-full-packet}

This calculation concerns the original polynomial at its original
marked parameter \(t=0\). It does not replace that polynomial by a
single-primary polynomial. All its primary pieces occur simultaneously
in the target of a proved map. The completion used below is in the
deformation variable \(u\); its coefficient, differential, and marked
fibre maps are specified explicitly.

\subsection{The fixed polynomial and the two complete complexes}

Fix the ordered nonempty finite set \(Z=(\rho_1,\ldots,\rho_r)\) of selected
actual zeros of \(g=2\xi\), with their complete orders \(m_\rho\ge1\).
Put
\[
 h(s)=\prod_{\rho\in Z}(s-\rho)^{m_\rho}
     =s^d+\sum_{b=0}^{d-1}h_bs^b,
 \qquad d=\sum_\rho m_\rho,
 \qquad
 \Phi(s)=\frac{s^{d+1}}{d+1}
       +\sum_{b=0}^{d-1}\frac{h_bs^{b+1}}{b+1}.
 \tag{FPB.1}\label{GraphV2:eq:fpb-fixed}
\]
Thus \(\Phi'=h\) and \(\Phi(0)=0\). In particular the constants
\(c_\rho:=\Phi(\rho)\) are fixed by the original polynomial; they
are not the constants of separately reassigned primitives.

Let \(A=\mathbb{C}[[u]]\), \(K=\mathbb{C}((u))\), and define
\[
 R=\mathbb{C}[s][[u]]
   =\varprojlim_n\mathbb{C}[s,u]/(u^n),
 \qquad
 \widehat{\mathcal C}=[R\xrightarrow{L}R\,ds],
 \qquad Lf=(u\partial_s f+h f)\,ds.
 \tag{FPB.2}\label{GraphV2:eq:fpb-global-complex}
\]
Degrees are zero and one, respectively. A coefficient of each power of
\(u\) is a polynomial in \(s\), with no common degree bound required.
Thus \(R\) is not being identified with \(A[s]\). It is complete and
separated for the \(u\)-adic topology. Its differential is continuous
for that topology.

For each \(\rho\), write \(y=s-\rho\), \(m=m_\rho\), and \(N=m+1\).
Define the entire polynomial unit germ
\[
 a_\rho(y)=\frac{N(\Phi(\rho+y)-\Phi(\rho))}{y^N},
 \qquad
 a_\rho(0)=\prod_{\sigma\ne\rho}(\rho-\sigma)^{m_\sigma}\ne0.
 \tag{FPB.3}\label{GraphV2:eq:fpb-a}
\]
The first equality is polynomial because \(\Phi(\rho+y)-\Phi(\rho)\)
vanishes to order \(N\). For the second, integrate
\(h(\rho+y)=y^m\prod_{\sigma\ne\rho}(\rho-\sigma+y)^{m_\sigma}\)
coefficient by coefficient. Choose and retain a specified number
\(b_\rho\) with \(b_\rho^N=a_\rho(0)\). The branch
\[
 A_\rho(y)=b_\rho
 \sum_{j\ge0}\binom{1/N}{j}
       \left(\frac{a_\rho(y)}{a_\rho(0)}-1\right)^j,
 \qquad w=yA_\rho(y),\qquad y=\psi_\rho(w)
 \tag{FPB.4}\label{GraphV2:eq:fpb-coordinate}
\]
is a formal coordinate with a uniquely specified formal inverse.
Indeed its linear coefficient is \(b_\rho\ne0\); equating the
coefficient of \(w^n\) in \(\psi A_\rho(\psi)=w\) determines the
coefficient of \(w^n\) in \(\psi\) uniquely by division by
\(b_\rho\), after the preceding coefficients have been determined.
The binomial identity gives \(A_\rho^N=a_\rho\). Consequently
\[
 \Phi(\rho+\psi_\rho(w))=c_\rho+\frac{w^N}{N},
 \qquad h(\rho+\psi_\rho(w))\psi_\rho'(w)=w^m.
 \tag{FPB.5}\label{GraphV2:eq:fpb-phase}
\]
These are equalities of formal series retaining the original constant.

Put \(R_\rho=\mathbb{C}[[w]][[u]]\), equipped with its \(u\)-adic
topology and the coefficientwise formal derivative in \(w\). Although
its elements can also be written as double formal series, only the
stated topology is used in the completion arguments. Set
\[
 \mathcal C_\rho=[R_\rho\xrightarrow{L_\rho}R_\rho\,dw],
 \qquad L_\rho f=(u\partial_w f+w^m f)\,dw.
 \tag{FPB.6}\label{GraphV2:eq:fpb-local-complex}
\]
The actual map of complexes is
\[
 F^0_\rho f=f(\rho+\psi_\rho(w)),\qquad
 F^1_\rho(f\,ds)=f(\rho+\psi_\rho(w))\psi_\rho'(w)\,dw,
 \qquad F=\bigoplus_{\rho\in Z}F_\rho.
 \tag{FPB.7}\label{GraphV2:eq:fpb-chain-map}
\]
Substitution is performed on each polynomial coefficient of \(u\).
The chain rule and (FPB.5) prove \(F^1_\rho L=L_\rho F^0_\rho\).
The Jacobian in degree one is part of this identity. The map \(F\)
is not asserted to be surjective on cochains.

\subsection{Complete remainder, inverse, and chain homotopies}

For \(f=\sum_{n\ge0}u^nf_n(s)\in R\), recursively divide
\[
 f_n-\partial_s q_{n-1}=h q_n+r_n,
 \qquad q_{-1}=0,\qquad \deg r_n<d.
 \tag{FPB.8}\label{GraphV2:eq:fpb-global-division}
\]
Each division is the original monic polynomial division. Then
\(f=Lq+r\), where \(q=\sum u^nq_n\), \(r=\sum u^nr_n\), and the
one-form symbol is understood. Conversely this identity determines
every coefficient by (FPB.8). In particular \(L\) is injective, and
\[
 R=L(R)\oplus\mathbb{C}[s]_{<d}[[u]],\qquad
 H^0(\widehat{\mathcal C})=0,\qquad
 \mathcal H:=H^1(\widehat{\mathcal C})
       =\bigoplus_{b=0}^{d-1}A[s^bds].
 \tag{FPB.9}\label{GraphV2:eq:fpb-global-remainder}
\]
For injectivity one can also take the first nonzero coefficient of
\(u\) in \(Lq=0\), obtaining \(h q_n=0\). Since
\(\mathbb{C}[s]\) is a domain, this is impossible.

In the local ring, divide every coefficient of
\(f_n-\partial_wq_{n-1}\in\mathbb{C}[[w]]\) uniquely by \(w^m\):
\[
 f_n-\partial_wq_{n-1}=w^m q_n+r_n,
 \qquad \deg r_n<m.
 \tag{FPB.10}\label{GraphV2:eq:fpb-local-division}
\]
The same coefficient recursion proves
\[
 H^0(\mathcal C_\rho)=0,
 \qquad \mathcal H_\rho:=H^1(\mathcal C_\rho)
       =\bigoplus_{a=0}^{m_\rho-1}A[w^a dw].
 \tag{FPB.11}\label{GraphV2:eq:fpb-local-remainder}
\]
No convergence of an infinite sum of fixed \(u\)-order is needed:
each recurrence is a single formal division in \(w\). The resulting
quotient and remainder maps, denoted \(Q,R_h\) globally and
\(Q_\rho,R_\rho^{\rm rem}\) locally, are continuous and \(A\)-linear.
Their constructions agree with the polynomial-exponential remainder
when the input is a polynomial, because that remainder has the same
unique decomposition. Thus (FPB.9) is also the completion of the
original global lattice, rather than a differently presented quotient.

In local coordinates the remainder has an explicit coefficient
formula. Define
\[
 p_{a,q}=\prod_{\ell=0}^{q-1}(a+1+\ell N),\qquad p_{a,0}=1.
\]
For \(0\le a<m\) and \(q\ge0\),
\[
 R_\rho^{\rm rem}(w^{a+qN})=(-u)^q p_{a,q}w^a,
 \qquad
 R_\rho^{\rm rem}(w^{m+qN})=0.
 \tag{FPB.12}\label{GraphV2:eq:fpb-monomial-remainder}
\]
To prove this, use
\(L_\rho(w^{j-m})=w^j+u(j-m)w^{j-N}\) for \(j\ge m\).
The case \(j=m\) gives zero remainder. Repeatedly reducing a monomial
in the other congruence classes gives precisely the displayed product.
For an infinite series in \(w\), its coefficient of \(u^q w^a\)
uses just the coefficient of \(w^{a+qN}\). For a series already
containing \(u\), the coefficient of any specified power of \(u\)
uses a finite sum. This proves (FPB.12) at the stated complete types.

Write the matrix induced by \(F\) in the original ascending monomial
columns and the ordered local rows \((\rho,a)\) as \(T(u)\). Then
\[
 \begin{split}
 T_{(\rho,a),b}(u)
 &=\sum_{q\ge0}(-u)^q p_{a,q}
   [y^{a+qN}](\rho+y)^b
       a_\rho(y)^{-(a+qN+1)/N},\\
 &\hspace{25mm}0\le b<d,\quad 0\le a<m_\rho.
 \end{split}
 \tag{FPB.13}\label{GraphV2:eq:fpb-T}
\]
Every fractional power uses the branch in (FPB.4), so its constant
is \(b_\rho^{-(a+qN+1)}\). For proof of this formula, formal
change of variable in a residue gives, for any formal germ
\(f(\rho+y)\in\mathbb{C}[[y]]\),
\[
 [w^j]f(\rho+\psi(w))\psi'(w)
   =[y^j]f(\rho+y)a_\rho(y)^{-(j+1)/N}.
 \tag{FPB.14}\label{GraphV2:eq:fpb-residue}
\]
Indeed the left side is the coefficient of \(w^{-1}dw\) in
\(f(\rho+\psi)\psi' w^{-j-1}dw\); substituting
\(w=yA_\rho(y)\) gives
\(f(\rho+y)y^{-j-1}A_\rho(y)^{-j-1}dy\). Formal residues are
unchanged by this substitution: for monomials other than exponent
\(-1\), this follows by differentiating a formal Laurent primitive;
for exponent \(-1\), \(w'/w=1/y+A_\rho'/A_\rho\) has residue one.
Combining (FPB.14) and (FPB.12) proves (FPB.13).

At \(u=0\), let \(E_h=\mathbb{C}[s]/h\) and
\(E_\rho^w=\mathbb{C}[w]/w^{m_\rho}\). Then
\[
 T_0:E_h\,ds\longrightarrow\bigoplus_\rho E_\rho^w\,dw,
 \qquad [f]\longmapsto
   ([f(\rho+\psi_\rho)\psi_\rho']\,dw)_\rho
 \tag{FPB.15}\label{GraphV2:eq:fpb-T0}
\]
is invertible. Here is its explicit inverse. Given local remainders
\(q_\rho(w)\), set
\[
 f_\rho(y)=q_\rho(yA_\rho(y))
       \frac{d}{dy}(yA_\rho(y))\pmod {y^{m_\rho}}.
 \tag{FPB.16}\label{GraphV2:eq:fpb-inverse-jet}
\]
Set \(H_\rho=h/(s-\rho)^{m_\rho}\) and let \(B_\rho(y)\) be the
truncated inverse of \(H_\rho(\rho+y)\) modulo \(y^{m_\rho}\),
calculated by the recursion for the inverse of a series with nonzero
constant. The original Hermite idempotent is the degree-below-\(d\)
remainder
\(e_\rho=\operatorname{rem}_h(H_\rho B_\rho(s-\rho))\).
The inverse of (FPB.15) is
\[
 (q_\rho)_\rho\longmapsto
  \operatorname{rem}_h\sum_\rho e_\rho(s)f_\rho(s-\rho).
 \tag{FPB.17}\label{GraphV2:eq:fpb-CRT}
\]
Each \(e_\rho\) is one modulo its own primary factor and zero
modulo every other one. The product of the distinct primary factors
divides a polynomial divisible by each: repeatedly apply the Bezout
identity to pairs of coprime factors. Thus their simultaneous jets
determine a unique remainder modulo \(h\). Finally
\(\psi'(w)w'(\psi(w))=1\) verifies (FPB.16) in both directions.
This proves (FPB.17), without assuming the desired comparison.

With \(T(u)=\sum_{n\ge0}u^nT_n\), define
\[
 V_0=T_0^{-1},\qquad
 V_n=-T_0^{-1}\sum_{i=1}^n T_iV_{n-i}\quad(n\ge1).
 \tag{FPB.18}\label{GraphV2:eq:fpb-Tinverse}
\]
For \(V(u)=\sum_{n\ge0}u^nV_n\), it follows coefficient by
coefficient that \(TV=1\). The same
recursion or the adjugate formula over \(A\), since \(\det T_0\ne0\),
gives \(VT=1\). Therefore
\[
 F_*: \mathcal H\xrightarrow[\ T(u)\ ]{\ \sim\ }
       \bigoplus_\rho\mathcal H_\rho
 \tag{FPB.19}\label{GraphV2:eq:fpb-isomorphism}
\]
is the required isomorphism of complete original lattices.

There are explicit chain homotopies as well. Let \(i_h\) and
\(i_{\rm loc}\) include remainder representatives into degree one,
and let \(p_h,p_{\rm loc}\) denote their remainder projections,
with zero maps in degree zero. The quotient maps \(Q,Q_{\rm loc}\)
are degree-minus-one maps and satisfy
\[
 1-i_hp_h=dQ+Qd,\qquad
 1-i_{\rm loc}p_{\rm loc}=dQ_{\rm loc}+Q_{\rm loc}d.
 \tag{FPB.20}\label{GraphV2:eq:fpb-contractions}
\]
Define \(J=i_hT^{-1}p_{\rm loc}\) as a map from the direct sum of
local complexes to the global complex. Since remainder vanishes on
every differential, this is a chain map. The identity
\(p_{\rm loc}F=Tp_h\) gives \(JF=i_hp_h\). Moreover
\(p_{\rm loc}(1-FJ)=0\), and hence
\[
 1-JF=dQ+Qd,
 \qquad
 1-FJ=d\bigl(Q_{\rm loc}(1-FJ)\bigr)
          +\bigl(Q_{\rm loc}(1-FJ)\bigr)d.
 \tag{FPB.21}\label{GraphV2:eq:fpb-homotopy}
\]
For the second identity, in degree one use zero remainder to write
\((1-FJ)f=L_{\rm loc}Q_{\rm loc}(1-FJ)f\); in degree zero,
\(FJ=0\) and \(Q_{\rm loc}L_{\rm loc}=1\). This also proves that
these formulas retain the cochain information through a specific
homotopy equivalence. The homotopies are \(A\)-linear; no assertion
that the chosen remainders are horizontal is needed.

\subsection{The exact connection and the constant determinant}

Localize the displayed complete complexes at \(u\). On the global
degrees zero and one respectively the original connection is
\[
 \nabla^0=\partial_u-\frac{\Phi(s)}{u^2}+\frac1u,
 \qquad
 \nabla^1=\partial_u-\frac{\Phi(s)}{u^2}.
 \tag{FPB.22}\label{GraphV2:eq:fpb-chain-connection}
\]
All operators are defined on \(R[u^{-1}]\): each element has a
finite lower bound on its powers of \(u\), and differentiation
preserves that property. Direct calculation gives
\([u\partial_s+h,\partial_u]=-\partial_s\) and
\([u\partial_s+h,-\Phi/u^2]=-h/u\), proving
\(\nabla^1L=L\nabla^0\). Under (FPB.7) the phase becomes exactly
\(c_\rho+w^N/N\). The change of variables and its Jacobian are
independent of \(u\), so \(F\) intertwines both degree connections
on these localized complexes.

In the unchanged global monomial basis, divide
\[
 \Phi(s)s^b=h(s)q_b(s)+r_b(s),\qquad \deg r_b<d,
 \quad 0\le b<d.
\]
Let \(C\) have columns \(r_b\) and \(B\) have columns
\(q_b'\). The derivative already has degree below \(d\).
Subtracting \(Lq_b\) gives the original connection matrix
\[
 G(u)=-C/u^2+B/u.
 \tag{FPB.23}\label{GraphV2:eq:fpb-G}
\]
For the local column \(w^a\),
\((c_\rho+w^N/N)w^a=w^m w^{a+1}/N+c_\rho w^a\).
The same reduction gives
\[
 D(u)=\bigoplus_{\rho\in Z}
  \operatorname{diag}_{0\le a<m_\rho}
    \left(-\frac{c_\rho}{u^2}
               +\frac{a+1}{N_\rho u}\right),
 \qquad T'+DT=TG.
 \tag{FPB.24}\label{GraphV2:eq:fpb-gauge}
\]
The second identity is equality of actual matrices over \(K\),
obtained by applying the proved chain intertwining to each original
basis column. It is also an all-order identity satisfied by the
coefficient formula (FPB.13).

The determinant is explicit and independent of \(u\):
\[
 \boxed{\displaystyle
 \det T(u)=\det T_0
   =\left(\prod_{\rho\in Z}
       b_\rho^{-m_\rho(m_\rho+1)/2}\right)
      \prod_{i<j}(\rho_j-\rho_i)^{m_{\rho_i}m_{\rho_j}}.}
 \tag{FPB.25}\label{GraphV2:eq:fpb-det}
\]
First consider \(T_0\). The ordinary jet matrix has entries
\(\binom ba\rho^{b-a}\). Its determinant is the confluent
Vandermonde product in (FPB.25), with coefficient one and the displayed
order. One elementary verification starts with the Vandermonde
determinant at distinct variables \(x_{i,a}\), whose determinant is
\(\prod_{(i,a)<(j,b)}(x_{j,b}-x_{i,a})\): it is alternating, divisible
by every displayed difference, and has that leading coefficient by
the identity permutation in the determinant. Apply successive divided
differences within each ordered cluster and then set every variable
in cluster \(i\) to \(\rho_i\). The divided-difference row of order
\(a\) becomes the Taylor-coefficient row \(f^{(a)}(\rho_i)/a!\).
The denominators cancel precisely the within-cluster Vandermonde
factors; each cross-cluster difference becomes \(\rho_j-\rho_i\).
This proves the asserted coefficient and sign. The map of a local
jet \(y^a dy\) to \(\psi(w)^a\psi'(w)dw\) is triangular with
diagonal \(b_\rho^{-(a+1)}\), proving the formula for \(\det T_0\).

For independence of \(u\), \(C\) is multiplication by \(\Phi\)
modulo \(h\). On the \(\rho\)-primary quotient,
\(\Phi(\rho+y)=c_\rho+O(y^{m_\rho+1})\), so it is exactly scalar
\(c_\rho\), with multiplicity \(m_\rho\). Hence
\(\operatorname{Tr}C=\sum_\rho m_\rho c_\rho\).
The quotient \(q_b\) has degree \(b+1\) and leading coefficient
\(1/(d+1)\), so \(B\) is upper triangular with diagonal
\((b+1)/(d+1)\) and trace \(d/2\). Also
\(\sum_{a=0}^{m-1}(a+1)/(m+1)=m/2\). Thus
\(\operatorname{Tr}G=\operatorname{Tr}D\).
Taking the trace of (FPB.24) after multiplication by \(T^{-1}\)
gives \((\det T)'=0\). In characteristic zero a formal series with
zero derivative is constant, proving (FPB.25).

Replacing the chosen branch by \(b_\rho\epsilon\),
\(\epsilon^{N_\rho}=1\), replaces \(w\) by \(\epsilon w\).
Since \(\epsilon^{a+qN+1}=\epsilon^{a+1}\), it replaces the local
rows of \(T\) by exactly
\(\operatorname{diag}(\epsilon^{-(a+1)})T_\rho\).
This is the complete branch-change map; the branch is not suppressed.

\subsection{The original arithmetic source and its Jacobian}

Keep the original \(\mathscr B,V,D=-x\partial_x,\Theta\), seed
\(\phi_*\), and Euler inverse \(F_h\), with
\(\mathcal M\Theta\phi_*=g\) and \(\mathcal M F_h=g/h\).
Explicitly, \(V\) consists of the even Schwartz functions on
\(\mathbb{R}\) satisfying \(\phi(0)=\int_{\mathbb{R}}\phi=0\),
and \(\mathscr B\) consists of smooth functions on
\(\mathbb{R}_{>0}\) with
\(\sup_{x>0}x^b|D^nF(x)|<\infty\) for every integer \(b\)
and every \(n\ge0\). The maps are
\(\Theta\phi(x)=2\sum_{n\ge1}\phi(nx)\) and
\(\mathcal MF(s)=\int_0^\infty F(x)x^s\,dx/x\);
\(\phi_*=(4\pi^2x^4-6\pi x^2)e^{-\pi x^2}\).
The original source maps are
\[
 \begin{gathered}
 \alpha_h(P)=P(D)\phi_*,\qquad
 \mathcal T_h(P)=P(D)F_h,\\
 \Theta\alpha_h(P)=\mathcal T_h(hP),\qquad
 J_h\mathcal T_h(P)=\upsilon_h[P],
 \qquad \upsilon_h=[g/h]\in E_h^\times.
 \end{gathered}
 \tag{FPB.26}\label{GraphV2:eq:fpb-theta}
\]
These are the inherited maps, with their original spaces and norms;
the coordinate calculation does not replace those spaces. In each
local primary factor retain every coefficient
\[
 v_\rho(y)=\sum_{j=0}^{m_\rho-1}
    \frac1{j!}\left(\frac{d^j}{ds^j}\frac{g(s)}{h(s)}
             \right)_{s=\rho}y^j\pmod {y^{m_\rho}}.
 \tag{FPB.27}\label{GraphV2:eq:fpb-unit}
\]
Its retained constant is exactly
\[
 v_\rho(0)=\frac{g^{(m_\rho)}(\rho)}
  {m_\rho!\,\prod_{\sigma\ne\rho}(\rho-\sigma)^{m_\sigma}}
  \ne0,
\]
because the selected order is complete. This equality follows by
factoring the leading Taylor term of \(g\) and the original factor
\((s-\rho)^{m_\rho}\) of \(h\).
The ordinary marked arithmetic operator is \(A_h=M_s\), and
the ordinary full-unit operator is \(U_h=M_{\upsilon_h}\).
Under (FPB.15), their exact marked images are
\[
 T_0A_hT_0^{-1}=\bigoplus_\rho M_{\rho+\psi_\rho(w)},\qquad
 T_0U_hT_0^{-1}=\bigoplus_\rho M_{v_\rho(\psi_\rho(w))}.
 \tag{FPB.28}\label{GraphV2:eq:fpb-marked-arithmetic}
\]
These multiplications are on \(E_\rho^w\,dw\), hence truncated
at the complete original order. To check them, compose the formulas
for \(T_0\) and its inverse; the two Jacobians cancel in a
conjugation, while the full Taylor series remains.

Let \(z_\rho=e_\rho(s)(s-\rho)^{m_\rho-1}\in E_h\). Then
\[
 T_0z_\rho=b_\rho^{-m_\rho}w^{m_\rho-1}dw
       \quad\hbox{in the \(\rho\)-factor, zero elsewhere},
 \qquad
 T_0U_hz_\rho=v_\rho(0)b_\rho^{-m_\rho}
                       w^{m_\rho-1}dw.
 \tag{FPB.29}\label{GraphV2:eq:fpb-terminal}
\]
Indeed \(\psi(w)=b_\rho^{-1}w+O(w^2)\) and
\(\psi'(w)=b_\rho^{-1}+O(w)\); all higher terms vanish at
degree \(m_\rho\). Both displayed terminal coefficients are nonzero.

Away from the marked fibre, extend the original matrices \(A_h,U_h\)
over \(A\) in the unchanged global monomial basis. Their exact local
operators are the explicitly computable matrices
\[
 \mathcal A(u)=T(u)A_hT(u)^{-1},\qquad
 \mathcal U(u)=T(u)U_hT(u)^{-1}.
 \tag{FPB.30}\label{GraphV2:eq:fpb-transported-operators}
\]
Equations (FPB.13) and (FPB.18) determine every coefficient of these
operators. They are not identified with naive multiplication on a
twisted de Rham quotient. The exact obstruction to that identification
is visible already on representatives: for a polynomial \(v\),
\[
 vLq=L(vq)-u v'q,\qquad
 R_h(vLq)=-uR_h(v'q).
 \tag{FPB.31}\label{GraphV2:eq:fpb-Leibniz}
\]
At \(u=0\) this defect vanishes. For an original representative
\(p\) of degree below \(d\), write \(vp=hq+r\), \(\deg r<d\).
Then the full deformed remainder is
\[
 R_h(vp)=r-uR_h(q'),
 \qquad
 T R_h(vp)=T r-uT R_h(q').
 \tag{FPB.32}\label{GraphV2:eq:fpb-mult-remainder}
\]
This proves the precise comparison with coefficientwise multiplication
by \([v]\) in \(E_h\), including its derivative term. The same
identity in the local complex has \(v'\) replaced by the derivative
of the actual pulled-back coefficient \(v(\rho+\psi(w))\).
For \(v=s\) and \(\deg p<d\), the quotient \(q\) is constant,
so the defect in (FPB.32) is zero on that section; (FPB.31) still
proves why this is not a multiplication map on arbitrary cochains.

\subsection{The precise attachment to the mixed support carriers}

The marked comparison contains two different, explicitly related
maps. The algebra map is
\[
 S_0:E_h\xrightarrow{\sim}\prod_\rho E_\rho^w,
 \qquad [f]\longmapsto([f(\rho+\psi_\rho)])_\rho.
 \tag{FPB.33}\label{GraphV2:eq:fpb-ring-map}
\]
It is a unital algebra isomorphism: the same CRT proof applies, with
the Jacobian omitted in both directions. Put
\(j=(\psi_\rho'(w))_\rho\), an invertible element of its target.
Then the actual one-form map is \(T_0=M_jS_0\), and it satisfies
\[
 T_0(a x)=S_0(a)T_0(x)\quad(a,x\in E_h).
 \tag{FPB.34}\label{GraphV2:eq:fpb-semilinear}
\]
Thus \((S_0,T_0)\) is an exact isomorphism of a coefficient algebra
and its rank-one module of marked one-forms. It is this pair that
attaches to the mixed carriers. In general \(T_0\) itself is not
a unital algebra map: it sends \(1\,ds\) to \(j\,dw\). If
multiplication by an invertible \(j\) were multiplicative, its value
on \(1\cdot1\) would force \(j=j^2\), hence \(j=1\).

For completeness, the split extension of a unital ring map \(f:R\to R'\)
is \(G(f)(\tau)=\tau\) and \(G(f)(r^\bullet)=f(r)^\bullet\).
The operations are the original ones: \(\tau\) is additive zero
and multiplicatively absorbing; supported addition and multiplication
are those of the ring. These formulas prove directly that \(G(f)\)
is a unital semiring map and that the supported zero
\(e=0^\bullet\) stays distinct from \(\tau\). In particular
\(G(S_0)\) retains the original absolute-base point and the supported
zero. The synchronized map
\[
 G\!\left(\prod_\rho E_\rho^w\right)
    \longrightarrow\prod_\rho G(E_\rho^w)
\]
sends \(\tau\) to the all-\(\tau\) tuple and a supported tuple to
the tuple with every coordinate supported. It does not introduce or
delete the other independent support masks in the target.

More generally every original independent coefficient-slot mask and
every original source-leg mask is retained as an external label. On
the module attached to a fixed coefficient face, apply \(F\), its
remainder maps, and its homotopies coefficientwise; on the label use
the identity. Inclusions, projections, and zero insertions of labelled
slots commute with this operation by linearity. A newly inserted
supported-zero slot remains supported-zero because its label is
unchanged, while a missing slot remains missing. The same statement
holds for a proper-source kernel carried as an additional labelled
summand: the map acts on its coefficient module and never quotients
out the summand. For multiplicative carrier arrows it is \(S_0\),
not the Jacobian module map, that supplies the coefficient-ring arrow.
On the deformed lattices the exact coefficient action is its original
action transported by \(T(u)\), as in (FPB.30); no unwarranted
multiplicativity of the remainder projection is used.

\subsection{All ordered tensor factors and formal character blocks}

For \(k\ge1\), use the original ordered complex on
\(\mathbb{C}[s_1,\ldots,s_k][[u]]\) with differential
\[
 \mathcal D=\sum_{i=1}^k
       (u\partial_{s_i}+h(s_i))\,ds_i\wedge(-).
 \tag{FPB.35}\label{GraphV2:eq:fpb-tensor-complex}
\]
This is the completed tensor product over \(A\) of the one-factor
complexes: modulo \(u^n\), the ordinary tensor product has polynomial
coefficients in all the \(s_i\); taking the inverse limit gives
precisely the displayed ring. The one-factor split homotopies extend
to these completed tensor products. On a tensor summand with a
contractible factor, apply its degree-minus-one homotopy with the
sign from all preceding degrees; the two cross terms cancel in the
tensor differential. Iterating leaves only the ordered top remainder
summand. Thus its cohomology is concentrated in degree \(k\), with
basis \(\prod_i s_i^{b_i}ds_1\wedge\cdots\wedge ds_k\),
\(0\le b_i<d\), and rank \(d^k\).

Apply (FPB.7) independently in every ordered factor. For clarity, the
local complex at a tuple \(\boldsymbol\rho\in Z^k\) is taken on
\(\mathbb{C}[[w_1,\ldots,w_k]][[u]]\), with its \(u\)-adic topology
and the differential \(\sum_i(u\partial_{w_i}+w_i^{m_{\rho_i}})
dw_i\wedge(-)\). This ring is not identified with the mere
\(u\)-adic completion of an algebraic tensor product of one-variable
formal-series rings. Its cohomology is computed by successive
one-variable versions of (FPB.10), with coefficient ring the formal
series in the remaining variables. Each division and derivative in
one variable commutes with those in another. The same split tensor
homotopies therefore leave the top remainder spanned over \(A\) by
\(\prod_iw_i^{a_i}\), \(0\le a_i<m_{\rho_i}\), and no other
cohomology. On an original monomial the separate reductions are the
product of the separate one-variable reductions. The target is the
direct sum over every tuple, and its top lattice map and phase are
exactly
\[
 \begin{gathered}
 T_k(u)=T(u)^{\otimes k},\\
 ds_1\wedge\cdots\wedge ds_k\longmapsto
  \prod_i\psi_{\rho_i}'(w_i)\,
       dw_1\wedge\cdots\wedge dw_k,\\
 \sum_i\Phi(s_i)=C_{\boldsymbol\rho}
                 +\sum_i\frac{w_i^{N_{\rho_i}}}{N_{\rho_i}},
 \quad C_{\boldsymbol\rho}=\sum_i c_{\rho_i}.
 \end{gathered}
 \tag{FPB.36}\label{GraphV2:eq:fpb-tensor-map}
\]
No permutation of the ordered one-forms was made, so no additional
sign occurs. The target connection on the column indexed by
\((\boldsymbol\rho,\boldsymbol a)\) is
\[
 \partial_u-\frac{C_{\boldsymbol\rho}}{u^2}
              +\frac{\beta_{\boldsymbol\rho,\boldsymbol a}}u,
 \qquad
 \beta_{\boldsymbol\rho,\boldsymbol a}
        =\sum_i\frac{a_i+1}{N_{\rho_i}},
 \quad 0\le a_i<m_{\rho_i}.
 \tag{FPB.37}\label{GraphV2:eq:fpb-tensor-connection}
\]
Tensoring (FPB.24), or differentiating \(T^{\otimes k}\) in each
factor, proves the exact intertwining with the original tensor-sum
connection. In degree \(j\), the original chain connection retains
the correction \((k-j)/u\). All tuples with the same value of
\(C_{\boldsymbol\rho}\) remain as distinct summands. Their direct
sum is the full block with that exponential constant; no assertion
that these sums are distinct enters the proof.

The original arithmetic tensor action is \(A^{(k)}=\sum_i A_{h,i}\),
with ordered full unit \(U_h^{\otimes k}\). Under the complete map
these become \(T_k A^{(k)}T_k^{-1}\) and
\(T_kU_h^{\otimes k}T_k^{-1}\). At \(u=0\), the images are the
actual multiplications by \(\sum_i(\rho_i+\psi_{\rho_i}(w_i))\)
and \(\prod_i v_{\rho_i}(\psi_{\rho_i}(w_i))\), respectively.
In particular
\[
 T_{k,0}U_h^{\otimes k}\bigotimes_i z_{\rho_i}
   =\left(\prod_i v_{\rho_i}(0)b_{\rho_i}^{-m_{\rho_i}}\right)
        \prod_i w_i^{m_{\rho_i}-1}
                  dw_1\wedge\cdots\wedge dw_k\ne0.
 \tag{FPB.38}\label{GraphV2:eq:fpb-tensor-terminal}
\]
For the original cyclic terminal in a tuple, put
\(K_{\boldsymbol\rho}=\sum_i(m_{\rho_i}-1)\).
The multinomial formula gives, in that actual primary corner,
\[
 \left(\sum_i(s_i-\rho_i)\right)^{K_{\boldsymbol\rho}}
   =\frac{K_{\boldsymbol\rho}!}{\prod_i(m_{\rho_i}-1)!}
                  \prod_i(s_i-\rho_i)^{m_{\rho_i}-1}.
 \tag{FPB.39}\label{GraphV2:eq:fpb-cyclic-terminal}
\]
Every other monomial has an exponent exceeding one of the retained
orders and vanishes. Its theta-weighted image under \(T_{k,0}\)
therefore has the coefficient in (FPB.38) multiplied by exactly the
factorial ratio in (FPB.39). This coefficient is nonzero over the
original complex numbers. No factorial, Taylor unit, or Jacobian has
been absorbed into a new choice of terminal vector.

Here are the precise formal monodromy conventions. For one local
factor put \(\zeta_\rho=e^{2\pi i/N_\rho}\),
\(\beta_{\rho,a}=(a+1)/N_\rho\), and
\[
 W_{\rho,j a}=\zeta_\rho^{j(a+1)}-1\quad(1\le j\le m_\rho),
 \qquad
 d_{\rho,a}=N_\rho^{\beta_{\rho,a}-1}
           e^{\pi i\beta_{\rho,a}}\Gamma(\beta_{\rho,a}).
 \tag{FPB.40}\label{GraphV2:eq:fpb-local-constants}
\]
These are the original homogeneous-model ordered contour constants:
the ray of argument \((\arg u+\pi)/N_\rho+2\pi j/N_\rho\)
gives \(d_{\rho,a}u^{\beta_{\rho,a}}\zeta_\rho^{j(a+1)}\)
after the positive radial substitution \(r^{N_\rho}/(N_\rho|u|)\),
and subtraction of the zeroth ray gives \(W_\rho\). The columns
are independent because a polynomial
\(\sum_{a=0}^{m-1}x_a(z^{a+1}-1)\) vanishing at all nontrivial
\(N\)-th roots also vanishes at \(1\) and has degree below \(N\),
so is zero. The gamma factors are nonzero since \(0<\beta<1\).

Adjoin formal symbols \(e^{c_\rho/u}\) and the fractional powers
\(u^{\beta_{\rho,a}}\), with the derivative specified by their
displayed exponents. A full formal period matrix for the original
global connection is
\[
 \Pi^{\rm form}(u)=
   \left(\bigoplus_\rho
       e^{c_\rho/u}W_\rho
            \operatorname{diag}_a(d_{\rho,a}u^{\beta_{\rho,a}})
   \right)T(u).
 \tag{FPB.41}\label{GraphV2:eq:fpb-formal-period}
\]
Equations (FPB.24) and (FPB.40) verify
\((\Pi^{\rm form})'=\Pi^{\rm form}G\) directly. Its exact
determinant, in the specified block and monomial orders, is
\[
 \det\Pi^{\rm form}
   =\det T_0\left(\prod_\rho\det W_\rho
                              \prod_{a=0}^{m_\rho-1}d_{\rho,a}\right)
          u^{d/2}\exp\left(\frac{\sum_\rho m_\rho c_\rho}{u}\right).
 \tag{FPB.42}\label{GraphV2:eq:fpb-period-det}
\]
Formal continuation \(\log u\mapsto\log u+2\pi i\) fixes the
integer-power matrix \(T(u)\) and the exponential symbols. It
therefore has period monodromy
\[
 M^{\rm form}=\bigoplus_\rho
    W_\rho\operatorname{diag}_a(\zeta_\rho^{a+1})W_\rho^{-1}.
 \tag{FPB.43}\label{GraphV2:eq:fpb-formal-monodromy}
\]
For the differential module's horizontal coefficient columns the
solutions of \((\partial_u+D)v=0\) instead have factors
\(e^{-c_\rho/u}u^{-\beta_{\rho,a}}\). Thus their monodromy
characters are the inverses of the period characters in (FPB.43).
This records the dual convention explicitly.

The ordered tensor period matrix is
\((\Pi^{\rm form})^{\otimes k}\), and its period character on
\((\boldsymbol\rho,\boldsymbol a)\) is
\(\exp(2\pi i\beta_{\boldsymbol\rho,\boldsymbol a})\).
These matrices are diagonalizable. On the formal regular factors the
unipotent logarithm is zero, and the dimension of the character-one
space inside the block with constant \(C\) is exactly
\[
 \#\left\{(\boldsymbol\rho,\boldsymbol a):
    \sum_i c_{\rho_i}=C,\quad
    \sum_i\frac{a_i+1}{N_{\rho_i}}\in\mathbb Z,\quad
    0\le a_i<m_{\rho_i}\right\}.
 \tag{FPB.44}\label{GraphV2:eq:fpb-formal-invariants}
\]
Repeated equal characters contribute their full multiplicity because
all their ordered summands were retained. Formula (FPB.44) is about
the formal regular character factors of this proved decomposition;
it does not identify actual analytic inertia invariants or finite-field
boundary stalks.

\subsection{Horizontal formal projectors and the unchanged source form}

Let \(P_\rho\) be the constant block projector on the \(\rho\)-summand
of the local direct sum. The actual formal projector in the original
global lattice is
\[
 \mathcal E_\rho(u)=T(u)^{-1}P_\rho T(u),\qquad
 \mathcal E_\rho(0)=M_{e_\rho},\qquad
 \mathcal E_\rho'+[G,\mathcal E_\rho]=0.
 \tag{FPB.45}\label{GraphV2:eq:fpb-projectors}
\]
The first identity defines all coefficients by (FPB.13) and
(FPB.18). The marked identity follows from the weighted CRT map,
whose Jacobian preserves each primary summand. Orthogonality,
idempotence, and sum one follow by conjugating the corresponding
identities for \(P_\rho\). Here algebraic orthogonality means
\(\mathcal E_\rho\mathcal E_\sigma=0\) when \(\rho\ne\sigma\);
it makes no assertion of orthogonality in the original source metric.
Differentiating the definition and using
\(T'=TG-DT\), with \([D,P_\rho]=0\), gives the displayed horizontal
identity. Thus the retained primary labels have explicitly horizontal
projectors over the original complete lattice, including when critical
values coincide; they have not merely been recovered from a dimension
count.

There is also a marked constraint on equal-critical-value off-diagonal
blocks. If \(\rho\ne\sigma\) and \(c_\rho=c_\sigma=c\), then
\[
 M_{e_\sigma}B M_{e_\rho}=0\quad\hbox{on }E_h.
 \tag{FPB.46}\label{GraphV2:eq:fpb-equal-c}
\]
Take the degree-below-\(d\) representative \(f\) of an element
supported on the \(\rho\)-factor. Multiplication by \(\Phi\)
modulo \(h\) is \(c f\), so the original quotient defining \(Bf\)
is exactly \(q=(\Phi-c)f/h\), and \(Bf=q'\). At \(\sigma\),
\(f\) vanishes to order at least \(m_\sigma\), while
\(\Phi-c\) vanishes to order at least \(m_\sigma+1\). Therefore
\(q\) vanishes to order at least \(m_\sigma+1\), and \(q'\)
vanishes to order at least \(m_\sigma\), proving (FPB.46).
This proves the exact compressed first-order relation without assuming
that all original critical values are different.

The original source form on the polynomial coefficient section is
\[
 (G_\theta)_{ab}=\frac1{2\pi}\int_{\mathbb{R}}
   \overline{(1/2+it)}^{\,a}(1/2+it)^b
       \left|\frac{g(1/2+it)}{h(1/2+it)}\right|^2dt,
 \quad 0\le a,b<d.
 \tag{FPB.47}\label{GraphV2:eq:fpb-original-metric}
\]
At a removable zero this integrand uses the entire quotient \(g/h\).
To verify its source type, put \(x=e^z\). The original
\(L^2(dx)\) norm of \(\mathcal T_h(P)\) becomes the
\(L^2(dz)\) norm of \(e^{z/2}\mathcal T_h(P)(e^z)\).
Its Fourier transform with the positive exponential convention is
\(P(1/2+it)(g/h)(1/2+it)\). Fourier Plancherel, with inverse
factor \(1/(2\pi)\), gives (FPB.47). The original
\(\mathscr B\) bounds imply square integrability and decay of all
needed derivatives, so the displayed integrals are finite. If the
norm is zero, the continuous original source function is zero;
its Mellin transform \(P(g/h)\) is zero. Since \(g/h\) is a
nonzero entire function, a polynomial \(P\) with that property is
zero. Hence \(G_\theta\) is positive definite.

On a phase-coordinate vector representing the polynomial section
\(q=T_0p\), the very same form is
\[
 G_{\rm phase}=T_0^{-*}G_\theta T_0^{-1}.
 \tag{FPB.48}\label{GraphV2:eq:fpb-pulled-metric}
\]
On the actual theta jet \(q_\theta=T_0U_hp\), the same source
form instead has matrix
\[
 G_{\rm theta\ jet}=(T_0U_h)^{-*}G_\theta(T_0U_h)^{-1}.
 \tag{FPB.49}\label{GraphV2:eq:fpb-theta-jet-metric}
\]
Here \(X^{-*}=(X^{-1})^*\). Substitution in the two quadratic
forms proves both identities, including the full unit and every
Jacobian factor. Their ordered tensor versions use the \(k\)-fold
Kronecker products of the displayed maps and \(G_\theta\).
There is no assertion that a possibly divergent formal \(T(u)\)
has a positive analytic Gram matrix at a nonzero numerical \(u\).

The distinction is exact at the map level. Equations (FPB.7),
(FPB.19), and (FPB.24) give a chain homotopy equivalence of the
specified \(u\)-adic complexes and an isomorphism of their localized
formal differential modules. No convergence of \(T(u)\), no map
from its formal local model rays to the original global rapid-decay
cycles, and no Stokes transition matrices has been proved here.
Consequently (FPB.41) is the specified formal period solution, whose
relation to the original global connection is the proved identity
\((\Pi^{\rm form})'=\Pi^{\rm form}G\); it is not declared to be
an evaluated global analytic contour integral. The full original
theta operator is related to it by (FPB.26)--(FPB.32) and their
ordered tensors, rather than by an assumed Frobenius or weight
intertwiner. In particular no arithmetic purity or Riemann-hypothesis
conclusion is inserted into the boundary comparison.
